%=====================================================================
\section{The surface water circuit}
\label{sec:hyd:circuit}

The water circuit sits between the wellhead and the regeneration coil: a pump that fixes the
loop mass flow $\md$, a set of splits and mixing junctions that decide \emph{how much} water
sees the well, the store and the coil, and a lumped accumulator that moves heat in time rather
than only in space. The modelling hypothesis is deliberately blunt --- incompressible water,
one constant $c_{p,w}$ evaluated at a representative loop temperature, adiabatic instantaneous
junctions, pipe runs with neither capacitance nor heat loss. The cost is stated at the end of
this section; the benefit is that every junction collapses from an enthalpy balance to a
flow-weighted average of temperatures, which is what makes the coupled fixed point of
Section~\ref{sec:hyd:coil} cheap enough to run $\num{29220}$ time steps.

\subsection{Topology}
\label{sec:hyd:topology}

Starting at the pump discharge, at $T_{\mathrm{wi}}$: at the \textbf{E/F bifurcation} a
fraction $w_{v\mathrm{E}}$ of the loop flow goes \emph{down the well} and the remainder
$w_{v\mathrm{F}}=1-w_{v\mathrm{E}}$ \emph{bypasses} it at $T_{\mathrm{wi}}$. Branch E is
produced at $T_{\mathrm{ret}}$ and re-mixes with branch F to give $T_{\mathrm{wr}}$. At the
next fork, $w_{v\mathrm{A}}$ enters the \textbf{accumulator} and leaves at $T_{\mathrm{wa}}$
while $w_{v\mathrm{B}}=1-w_{v\mathrm{A}}$ heads straight for the coil; the accumulator outlet
is then split again, $w_{v\mathrm{C}}$ of it returning to the coil inlet and
$w_{v\mathrm{D}}=1-w_{v\mathrm{C}}$ bypassing the coil. The coil inlet blend is
$T_{\mathrm{wm}}$ and its water outlet $T_{\mathrm{wh}}$; mixing that with branch D gives the
loop return $T_{\mathrm{wo}}$, and the pump closes the loop back to $T_{\mathrm{wi}}$.

Table~\ref{tab:hyd:valves} collects the split fractions. In the baseline $w_{v\mathrm{E}}=1$,
so the F bypass is closed and the well carries the full loop flow; the branch exists because a
well-flow schedule is the natural optimisation variable, not because the baseline uses it.

\begin{table}[htbp]
\centering
\caption{Valve split fractions of the water circuit (Section~K of the
configuration cell). The complements $w_{v\mathrm{B}},w_{v\mathrm{D}},
w_{v\mathrm{F}}$ are not free parameters: each is $1$ minus its partner.}
\label{tab:hyd:valves}
\begin{tabular}{llS[table-format=1.2]l}
\toprule
Symbol & Meaning & {Baseline} & Complement \\
\midrule
$w_{v\mathrm{A}}$ & loop flow into the accumulator & 0.30 & $w_{v\mathrm{B}}=0.70$ \\
$w_{v\mathrm{C}}$ & accumulator outlet returned to the coil & 0.50 & $w_{v\mathrm{D}}=0.50$ \\
$w_{v\mathrm{E}}$ & loop flow sent down the well & 1.00 & $w_{v\mathrm{F}}=0$ \\
$w_{v\mathrm{E,min}}$ & floor imposed on the well flow & 0.05 & --- \\
\bottomrule
\end{tabular}
\end{table}

\subsection{The E/F bifurcation, and why branch F is not just a mixing node}
\label{sec:hyd:EF}

Write an enthalpy balance on the junction where the produced stream $w_{v\mathrm{E}}\md$ at
$T_{\mathrm{ret}}$ meets the bypass $w_{v\mathrm{F}}\md$ at $T_{\mathrm{wi}}$. With $c_{p,w}$
constant across the junction the specific heats cancel and only the flow weights survive:
\begin{equation}
\boxed{\;T_{\mathrm{wr}}=w_{v\mathrm{E}}\,T_{\mathrm{ret}}
+\left(1-w_{v\mathrm{E}}\right)T_{\mathrm{wi}}\;}
\label{eq:hyd:Twr}
\end{equation}
Setting $w_{v\mathrm{E}}=1$ returns $T_{\mathrm{wr}}=T_{\mathrm{ret}}$, so the bifurcation is
backward compatible with the plain single-well loop.

Equation~\eqref{eq:hyd:Twr} looks like a trivial dilution, and if the well were a fixed
temperature source it would be. It is not. The DBHE model is \emph{parametrised by the flow
through it}: the well is re-evaluated at $w_{v\mathrm{E}}\md$, which sets the annulus and
tubing velocities, hence the Reynolds numbers, hence the Gnielinski film coefficients, hence
the transfer units $\mathrm{NTU}_{at}$, $\mathrm{NTU}_{bw}$ and the effective wall conductance
$G_{\mathrm{eff}}$ that drives the formation. Closing branch F therefore does two opposing
things at once: it sends more water down the hole, \emph{increasing} the extracted duty, and it
shortens the residence time while changing the internal conductances, \emph{decreasing}
$T_{\mathrm{ret}}$. Branch F is a temperature-versus-duty knob on the subsurface, not a mixing
valve.

This is why the well model is split in two. The expensive, flow-independent objects --- the
GITT eigenmodes, the tail resistance, the persistent wall-temperature state --- are built once
at initialisation; only the cheap flow-dependent terms are refreshed each step, which makes a
time-varying $w_{v\mathrm{E}}$ affordable.

\paragraph{A consistency caveat.} The code refreshes the well terms at
$\max(w_{v\mathrm{E}},w_{v\mathrm{E,min}})$ with $w_{v\mathrm{E,min}}=\num{0.05}$ --- a
scheduled shutdown leaves a trickle circulating rather than a stagnant column --- while the
junction Eq.~\eqref{eq:hyd:Twr} uses the raw $w_{v\mathrm{E}}$. For any $w_{v\mathrm{E}}\ge
w_{v\mathrm{E,min}}$, including the baseline, the two coincide; below the floor they do not,
and a future scheduling study must either clamp the schedule or apply the floor in both places.

\subsection{Mixing junctions and the mass-conservation check}
\label{sec:hyd:mixing}

Two streams reach the coil inlet: the direct branch B, carrying $w_{v\mathrm{B}}\md$ at
$T_{\mathrm{wr}}$, and the recycled branch C, carrying $w_{v\mathrm{C}}w_{v\mathrm{A}}\md$ at
the accumulator outlet $T_{\mathrm{wa}}$. The fraction of the loop flow that actually passes
through the coil is therefore
\begin{equation}
\phi_{\mathrm{coil}}
=w_{v\mathrm{B}}+w_{v\mathrm{C}}\,w_{v\mathrm{A}}
=\left(1-w_{v\mathrm{A}}\right)+w_{v\mathrm{C}}\,w_{v\mathrm{A}} ,
\label{eq:hyd:fraccoil}
\end{equation}
and the same enthalpy-balance argument that gave Eq.~\eqref{eq:hyd:Twr} gives the coil inlet
and the loop return:
\begin{align}
T_{\mathrm{wm}}&=\frac{w_{v\mathrm{B}}\,T_{\mathrm{wr}}
+w_{v\mathrm{C}}w_{v\mathrm{A}}\,T_{\mathrm{wa}}}{\phi_{\mathrm{coil}}},
\label{eq:hyd:Twm}\\[2pt]
T_{\mathrm{wo}}&=\phi_{\mathrm{coil}}\,T_{\mathrm{wh}}
+w_{v\mathrm{D}}w_{v\mathrm{A}}\,T_{\mathrm{wa}} .
\label{eq:hyd:Two}
\end{align}
Equation~\eqref{eq:hyd:Twm} is normalised by $\phi_{\mathrm{coil}}$ because it is an average
over the coil stream only; Eq.~\eqref{eq:hyd:Two} is not, because it is an average over the
whole loop. That last statement holds only if the weights close, so check it rather than assert
it. Substituting $w_{v\mathrm{B}}=1-w_{v\mathrm{A}}$ and $w_{v\mathrm{D}}=1-w_{v\mathrm{C}}$,
\begin{equation}
\phi_{\mathrm{coil}}+w_{v\mathrm{D}}w_{v\mathrm{A}}
=\left(1-w_{v\mathrm{A}}\right)+w_{v\mathrm{C}}w_{v\mathrm{A}}
+\left(1-w_{v\mathrm{C}}\right)w_{v\mathrm{A}}
=1 .
\label{eq:hyd:massclose}
\end{equation}
Equation~\eqref{eq:hyd:Two} is thus a proper convex combination for \emph{any} valve setting,
and the circuit conserves mass identically. This is not decoration: because the loop is closed
and iterated to a fixed point, a weight set that did not close would inject a spurious enthalpy
source that the iteration would happily converge to, and the error would look like a physical
result.

\paragraph{The coil throttle is currently a constant.} It is tempting to read
$\phi_{\mathrm{coil}}$ as the plant's heat-delivery actuator, and eventually it will be. In the
code of record it is not: $\phi_{\mathrm{coil}}$ is evaluated \emph{once}, before the time
loop, from the fixed constants of Table~\ref{tab:hyd:valves}, so with
$w_{v\mathrm{A}}=\num{0.30}$ and $w_{v\mathrm{C}}=\num{0.50}$,
\begin{equation}
\phi_{\mathrm{coil}}=0.70+0.50\times0.30=\num{0.85},
\qquad
w_{v\mathrm{D}}w_{v\mathrm{A}}=\num{0.15},
\label{eq:hyd:fracnum}
\end{equation}
for the whole \num{20}-year run. The humidity actuator in the present model is the regeneration
vent damper, not the water throttle. Making $\phi_{\mathrm{coil}}$ a controlled variable ---
throttling water away from the coil when the wheel needs less regeneration and banking the
surplus in the accumulator --- is the intended next step, and
Eqs.~\eqref{eq:hyd:fraccoil}--\eqref{eq:hyd:massclose} already admit a time-varying
$w_{v\mathrm{A}},w_{v\mathrm{C}}$ without breaking the mass balance.

\subsection{The thermal accumulator}
\label{sec:hyd:acc}

The store --- a packed sand or concrete box, or a water tank --- is one lumped capacitance: a
single temperature $T_{\mathrm{acc}}$, no internal gradients, no stratification. Its heat
capacity is
\begin{equation}
C_{\mathrm{acc}}=\rho_{\mathrm{acc}}\,c_{p,\mathrm{acc}}\,V_{\mathrm{acc}}
=\num{1600}\times\num{830}\times\num{2000}
\approx\SI{2.66e9}{\joule\per\kelvin}.
\label{eq:hyd:Cacc}
\end{equation}
Water passing through it at $w_{v\mathrm{A}}\md$ does not reach $T_{\mathrm{acc}}$; the
approach is set by one heat-exchange effectiveness $\varepsilon_{\mathrm{acc}}$, defined as
always by actual over maximum possible temperature change, so the exit temperature is
\begin{equation}
T_{\mathrm{wa}}=T_{\mathrm{wr}}-\varepsilon_{\mathrm{acc}}
\left(T_{\mathrm{wr}}-T_{\mathrm{acc}}\right),
\qquad \varepsilon_{\mathrm{acc}}=\num{0.70}.
\label{eq:hyd:Twa}
\end{equation}
As $\varepsilon_{\mathrm{acc}}\to1$ the water leaves at the store temperature; as
$\varepsilon_{\mathrm{acc}}\to0$ the store is invisible. With no through-flow the code returns
$T_{\mathrm{wa}}=T_{\mathrm{acc}}$, keeping Eq.~\eqref{eq:hyd:Twm} finite in the degenerate
case.

An energy balance on the store --- what the water gives up minus what leaks to the surroundings
--- is the governing ODE:
\begin{equation}
\boxed{\;
C_{\mathrm{acc}}\frac{\mathrm{d}T_{\mathrm{acc}}}{\mathrm{d}t}
=\underbrace{\varepsilon_{\mathrm{acc}}\,w_{v\mathrm{A}}\md\,c_{p,w}
\left(T_{\mathrm{wr}}-T_{\mathrm{acc}}\right)}_{\dot Q_{\mathrm{charge}}}
-\underbrace{(\UA)_{\mathrm{acc}}\left(T_{\mathrm{acc}}-T_{\mathrm{amb}}\right)}
_{\dot Q_{\mathrm{loss}}}\;}
\label{eq:hyd:accode}
\end{equation}
The charge term is exactly the enthalpy drop of the through-flow implied by
Eq.~\eqref{eq:hyd:Twa}, since
\begin{equation*}
w_{v\mathrm{A}}\md\,c_{p,w}\left(T_{\mathrm{wr}}-T_{\mathrm{wa}}\right)
=\varepsilon_{\mathrm{acc}}\,w_{v\mathrm{A}}\md\,c_{p,w}
\left(T_{\mathrm{wr}}-T_{\mathrm{acc}}\right).
\end{equation*}
The store and the water agree on how much heat changed hands, so the model is energy consistent
at this junction --- not automatic when an effectiveness and a capacitance are written by
different hands.

The code integrates Eq.~\eqref{eq:hyd:accode} with a forward Euler step, committed once per
time step after the fixed point has converged:
\begin{equation}
T_{\mathrm{acc}}^{\,k+1}=T_{\mathrm{acc}}^{\,k}
+\frac{\Delta t}{C_{\mathrm{acc}}}
\Big[\varepsilon_{\mathrm{acc}}\,w_{v\mathrm{A}}\md\,c_{p,w}
\big(T_{\mathrm{wr}}^{\,k}-T_{\mathrm{acc}}^{\,k}\big)
-(\UA)_{\mathrm{acc}}\big(T_{\mathrm{acc}}^{\,k}-T_{\mathrm{amb}}\big)\Big],
\label{eq:hyd:accupdate}
\end{equation}
with $\Delta t=\SI{21600}{\second}$ (a quarter day). Explicit Euler on a linear relaxation is
stable only for $\Delta t\le C_{\mathrm{acc}}/
(\varepsilon_{\mathrm{acc}}w_{v\mathrm{A}}\md\,c_{p,w}+(\UA)_{\mathrm{acc}})$, which a store of
\SI{2.66e9}{\joule\per\kelvin} driven by a few kilograms per second satisfies by three orders
of magnitude. First-order accuracy is adequate because the store's own time constants are
months while the forcing is resolved four times a day.

The sign of $\dot Q_{\mathrm{charge}}$ takes care of itself: when
$T_{\mathrm{wr}}>T_{\mathrm{acc}}$ the store charges, and when the formation has been drawn
down over several years and $T_{\mathrm{wr}}$ falls below the store, the same term goes
negative and the store discharges into the loop. That reversal is the entire purpose of the
component, and it needs no switch.

\paragraph{The loss term is not a detail.} It is easy to treat $(\UA)_{\mathrm{acc}}$ as a
rounding error, because for a single overnight cycle it is one. With
$(\UA)_{\mathrm{acc}}=\SI{150}{\watt\per\kelvin}$ the store's loss time constant is
\begin{equation}
\tau_{\mathrm{loss}}=\frac{C_{\mathrm{acc}}}{(\UA)_{\mathrm{acc}}}
=\frac{\SI{2.66e9}{\joule\per\kelvin}}{\SI{150}{\watt\per\kelvin}}
\approx\SI{1.77e7}{\second}\approx\SI{205}{\day},
\label{eq:hyd:tauloss}
\end{equation}
so over one night an excess above ambient decays by about \SI{0.5}{\percent} --- negligible, as
expected. But the store is not proposed as an overnight buffer; it is proposed as the thing
that carries heat from a fresh year into a depleted one. On that horizon the same number reads
very differently: an unforced excess retains only $\exp(-365/205)\approx\SI{17}{\percent}$ of
itself after a year and loses about \SI{14}{\percent} in a single month. Whether "charge now,
recover later" is engineering or wishful thinking is decided entirely by $(\UA)_{\mathrm{acc}}$
integrated over the storage duration, and Eq.~\eqref{eq:hyd:tauloss} deserves a sensitivity
study before any claim about seasonal storage is made. Note also that the sink temperature is
bound to the outdoor air, $T_{\mathrm{amb}}=T_{\mathrm{oa}}$, so under weather-file forcing the
loss will track the season automatically --- largest in summer, which is precisely when the
store is being charged.

\subsection{Frictional pressure drop and pump power}
\label{sec:hyd:pump}

The coaxial well is two passages in series: down the annulus, up the tubing. Each contributes a
Darcy--Weisbach loss over the full depth $L$, and the model sums them:
\begin{equation}
\Delta p=\sum_{i\in\{a,t\}}f_i\,\frac{L}{D_{h,i}}\,
\frac{\rho\,v_i^{2}}{2},
\qquad
\mathrm{Re}_i=\frac{\rho\,|v_i|\,D_{h,i}}{\mu} .
\label{eq:hyd:dp}
\end{equation}
The friction factor is a two-branch correlation:
\begin{equation}
f_i=
\begin{cases}
\dfrac{64}{\mathrm{Re}_i}, & \mathrm{Re}_i<2300 \quad\text{(fully developed laminar)},\\[10pt]
\left(0.790\ln \mathrm{Re}_i-1.64\right)^{-2}, & \mathrm{Re}_i\ge2300
\quad\text{(smooth-pipe Petukhov fit)} .
\end{cases}
\label{eq:hyd:friction}
\end{equation}
The laminar branch is exact for a circular pipe; the turbulent branch is Petukhov's explicit
fit, chosen so that hydraulics and heat transfer stay consistent --- it is the same $f$ the
Gnielinski correlation uses for the film coefficients, so a single friction factor serves both.

\paragraph{What the switch costs.} The two branches do not meet. At $\mathrm{Re}=2300$ the
laminar expression gives $f=\num{0.0278}$ and the Petukhov fit gives $f=\num{0.0499}$ --- a
factor of \num{1.8} jump, with no blending applied --- and Petukhov's fit, correlated for
$\mathrm{Re}\gtrsim\num{e4}$, is being extrapolated down to \num{2300}. In the design regime
both passages are comfortably turbulent and neither issue bites; but a low-flow schedule,
exactly what a time-varying $w_{v\mathrm{E}}$ produces, will walk the annulus through the
discontinuity and $\Delta p$ will show a step that is numerical, not physical.

Hydraulic power divided by the wire-to-water efficiency gives the pump consumption, and the
dissipated power is returned to the water as a temperature rise on the loop return:
\begin{equation}
P_{\mathrm{pump}}=\frac{\Delta p\,\dot V}{\eta_{\mathrm{pump}}},
\qquad \eta_{\mathrm{pump}}=\num{0.65},
\qquad
T_{\mathrm{wi}}=T_{\mathrm{wo}}+\frac{P_{\mathrm{pump}}}{\md\,c_{p,w}} .
\label{eq:hyd:pump}
\end{equation}
Equation~\eqref{eq:hyd:pump} closes the loop, and with it the water circuit.

\paragraph{Scope note.} Three simplifications live in
Eqs.~\eqref{eq:hyd:dp}--\eqref{eq:hyd:pump} and should be named rather than buried. First,
friction is evaluated on the \emph{well} stream only: surface piping, fittings, and the
throttling loss required to push $w_{v\mathrm{E}}\md$ through a deep, high-resistance well
while a near-zero-resistance bypass sits in parallel, are not counted. Second, $\dot V$ in
Eq.~\eqref{eq:hyd:pump} is the \emph{well} volumetric flow while the temperature rise is spread
over the \emph{loop} flow $\md$; these coincide only at $w_{v\mathrm{E}}=1$, the baseline, so
any partial-flow study inherits an attribution error. Third, the pump's thermal effect is small
next to the coil duty; the term is retained so that the loop energy balance closes exactly, not
because it matters numerically.

%=====================================================================
\section{The regeneration coil}
\label{sec:hyd:coil}

The coil is the plant's water-to-air interface: geothermal water inside the tubes, regeneration
air across the fins, no phase change on either side. Because the air is heated and nothing
condenses, the process is purely sensible and the humidity ratio passes through unchanged,
$w_{\mathrm{ro}}=w_{\mathrm{ri}}$. Everything the coil does is captured by one duty and two
outlet temperatures.

\subsection{Why effectiveness--NTU and not LMTD}
\label{sec:hyd:whyntu}

Both methods describe the same exchanger; the choice is decided by what is known and by how
often the calculation must run. This is a \emph{rating} problem: the inlet temperatures, the
capacity rates and $\UA$ are
known, and the outlets are wanted. The LMTD method builds the duty from a log-mean temperature
difference computed from the \emph{outlets} --- the very quantities being sought --- so it must
be iterated, and each iteration takes a logarithm of a difference that passes through zero when
the streams pinch. The $\varepsilon$--$\mathrm{NTU}$ method inverts that dependency:
effectiveness is a closed-form function of $\mathrm{NTU}$ and the capacity ratio alone, both
known before the solve, so the coil is evaluated once and exactly.

That matters here more than usual, because the coil does not sit on its own. It is evaluated
inside the coupled fixed point on $(T_{\mathrm{wi}},T_{\mathrm{ri}},w_{\mathrm{ri}})$, which
runs up to \num{40} iterations; that fixed point is itself re-run inside the humidity
back-solve, which allows up to \num{12} outer iterations plus two anchoring solves; and the
whole thing repeats for \num{29220} time steps. A coil needing its own inner iteration would
put a third loop inside the second and a non-smooth convergence criterion inside a relaxed
fixed point --- a reliable way to stall the outer solve.

\subsection{Capacity rates and the effectiveness relation}
\label{sec:hyd:coileqs}

The hot side is the coil water stream, the cold side the regeneration air:
\begin{equation}
C_h=\phi_{\mathrm{coil}}\,\md\,c_{p,w},
\qquad
C_c=\mdr\,c_{p,a},
\qquad c_{p,a}=\SI{1006}{\joule\per\kilogram\per\kelvin}.
\label{eq:hyd:caprates}
\end{equation}
From these follow the three quantities the method needs:
\begin{equation}
C_{\min}=\min\left(C_h,C_c\right),\quad
C_{\max}=\max\left(C_h,C_c\right),\quad
C_r=\frac{C_{\min}}{C_{\max}},\quad
\mathrm{NTU}=\frac{\UA_{\mathrm{coil}}}{C_{\min}} .
\label{eq:hyd:ntudef}
\end{equation}
$C_{\min}$ is the stream that runs out of capacity first, and it alone sets the thermodynamic
ceiling on the duty: no exchanger, however large, can transfer more than
$C_{\min}(T_{\mathrm{wm}}-T_{\mathrm{ri}})$, since doing so would drive that stream past the
other's inlet temperature. $C_r$ measures how badly the streams are matched; $\mathrm{NTU}$
measures how much exchanger has been bought relative to the stream that limits it.

The coil is a finned-tube crossflow unit with \emph{both fluids unmixed} --- water confined to
tubes, air to fin passages, so neither stream can even out its own transverse temperature
profile. This is the arrangement the code selects by default and the one used throughout the
facility model. Its standard closed-form approximation is
\begin{equation}
\boxed{\;\varepsilon=1-\exp\!\left\{\frac{\mathrm{NTU}^{0.22}}{C_r}
\left[\exp\!\left(-C_r\,\mathrm{NTU}^{0.78}\right)-1\right]\right\}\;}
\label{eq:hyd:epscross}
\end{equation}
This is a correlation, not an exact solution --- the exact unmixed--unmixed crossflow result is
an infinite series --- accurate to about \SI{1}{\percent} over the usual design range. A
counterflow branch also exists in the code for comparison, but it is not the default and the
facility driver never selects it; quoting counterflow effectiveness here would overstate the
coil, since counterflow is the best possible arrangement and unmixed crossflow is not.

\subsection{Duty and outlets}
\label{sec:hyd:duty}

Effectiveness is actual duty over the thermodynamic maximum, so the duty and both outlets
follow immediately from an energy balance on each stream:
\begin{equation}
Q_{\mathrm{coil}}=\varepsilon\,C_{\min}
\left(T_{\mathrm{wm}}-T_{\mathrm{ri}}\right),
\qquad
T_{\mathrm{ro}}=T_{\mathrm{ri}}+\frac{Q_{\mathrm{coil}}}{C_c},
\qquad
T_{\mathrm{wh}}=T_{\mathrm{wm}}-\frac{Q_{\mathrm{coil}}}{C_h} .
\label{eq:hyd:coilout}
\end{equation}
$T_{\mathrm{ro}}$ is the regeneration temperature delivered to the desiccant wheel;
$T_{\mathrm{wh}}$ returns to the mixing node of Eq.~\eqref{eq:hyd:Two}. Since $\varepsilon<1$
strictly, $T_{\mathrm{ro}}<T_{\mathrm{wm}}$ always: the coil approaches the water inlet
temperature but never reaches it, and no amount of $\UA$ changes that once $C_c$ is not the
limiting stream.

Equation~\eqref{eq:hyd:coilout} carries no sign guard. If the regeneration air ever arrives
hotter than the coil water, $Q_{\mathrm{coil}}$ goes negative and the coil cools the air ---
physically correct for a sensible exchanger, and deliberately left in, but it makes a negative
$Q_{\mathrm{coil}}$ a diagnostic worth reading rather than an error.

\paragraph{One approximation worth naming.} The cold-side capacity rate uses the dry-air
specific heat $c_{p,a}=\SI{1006}{\joule\per\kilogram\per\kelvin}$ rather than the moist-air
value $c_{p,a}+1860\,w$. At a humidity ratio of \SI{20}{\gram\per\kilogram} this understates
$C_c$ by about \SI{3.7}{\percent}, which propagates into $\mathrm{NTU}$, $\varepsilon$ and
$T_{\mathrm{ro}}$. The load side of the model does carry the moist-air correction, so the
inconsistency is confined to the coil and is small; it is recorded here so that it is not
rediscovered as a bug.

\subsection{The temperature--throughput trade-off}
\label{sec:hyd:tradeoff}

This is the coil's controlling design tension and it constrains the whole plant, so it is worth
deriving rather than asserting. Take the air side first, at a fixed water stream. If the air
flow is small, $C_c\ll C_h$, then
$C_{\min}=C_c$ and $\mathrm{NTU}=\UA/C_c$ is large: the air leaves close to the water inlet,
$T_{\mathrm{ro}}\to T_{\mathrm{wm}}$, but the duty $Q=\varepsilon
C_c(T_{\mathrm{wm}}-T_{\mathrm{ri}})$ is small because $C_c$ is. \emph{Hot air, little heat.}
If the air flow is large, $C_c\gg C_h$, then $C_{\min}=C_h$, the duty saturates at $\varepsilon
C_h(T_{\mathrm{wm}}-T_{\mathrm{ri}})$ --- the water is now the bottleneck --- and the
temperature rise $Q/C_c$ tends to zero. \emph{Much air, cool air.} The wheel needs a minimum
$T_{\mathrm{ro}}$ to regenerate at all; the plant needs enough $Q$ to cover the latent load.
The two requirements pull in opposite directions and the air-side design sits on that frontier.

The water side needs a more careful statement than it is usually given. Taken in isolation, at
\emph{fixed} $T_{\mathrm{wm}}$, raising the water flow raises $C_h$ and raises both $Q$ and
$T_{\mathrm{ro}}$; there is no trade-off internal to the coil, only diminishing returns as the
duty saturates at the air-limited value $\varepsilon C_c(T_{\mathrm{wm}}-T_{\mathrm{ri}})$. The
trade-off appears only once the coil is closed on the well. $T_{\mathrm{wm}}$ is not
independent: through Eqs.~\eqref{eq:hyd:Twr} and \eqref{eq:hyd:Twm} it is inherited from
$T_{\mathrm{ret}}$, and the DBHE delivers less temperature the harder it is pumped. Raising the
loop flow therefore raises the well's duty and lowers $T_{\mathrm{ret}}$, and the coil passes
that on faithfully: \emph{more water raises the duty and lowers the temperature the wheel
sees}. The mechanism is the well, not the exchanger.

The same mechanism explains why one well starves a large plant. Scaling the hatchery scales
$\mdr$ and hence $C_c$, but the water stream is set by a single borehole. Once $C_h<C_c$ the
coil is water-limited, the duty is pinned at $\varepsilon
C_h(T_{\mathrm{wm}}-T_{\mathrm{ri}})$, and the delivered temperature rise $Q/C_c$ collapses
towards zero as $C_c$ grows. No amount of coil substitutes for a second well.

\subsection{Sizing rule for the conductance}
\label{sec:hyd:sizing}

Fixing $\UA_{\mathrm{coil}}$ as an absolute constant would silently mis-size the coil whenever
the plant is scaled --- a ten-million-egg hatchery served by an exchanger sized for one
million. The model therefore ties the conductance to the stream it serves:
\begin{equation}
\UA_{\mathrm{coil}}=(ua)_{\mathrm{coil}}\,\mdr ,
\qquad
(ua)_{\mathrm{coil}}=\SI{3000}{\watt\per\kelvin}\ \text{per}\
\si{\kilogram\per\second},
\label{eq:hyd:UAcoil}
\end{equation}
with $\mdr$ itself derived from the process air flow. The choice of \num{3000} is not
arbitrary: it fixes the number of transfer units against the air stream at
\begin{equation}
\mathrm{NTU}\big|_{C_{\min}=C_c}
=\frac{(ua)_{\mathrm{coil}}\,\mdr}{\mdr\,c_{p,a}}
=\frac{\num{3000}}{\num{1006}}\approx\num{2.98},
\label{eq:hyd:NTUair}
\end{equation}
independent of plant size --- a conventional finned-coil design point, giving
$\varepsilon\approx\num{0.95}$ in the small-$C_r$ limit. The scaling is therefore
self-consistent: enlarge the plant and the coil grows with it, keeping the air-side design
point fixed. What does \emph{not} stay fixed is $C_r$, which is exactly the starvation
mechanism of Section~\ref{sec:hyd:tradeoff}.

\paragraph{What Eq.~\eqref{eq:hyd:UAcoil} hides.} The conductance is physically three
resistances in series,
\begin{equation}
\frac{1}{\UA_{\mathrm{coil}}}
=\frac{1}{h_w A_w}+R_{\mathrm{wall}}+\frac{1}{\eta_o h_a A_a},
\label{eq:hyd:UAdecomp}
\end{equation}
with $h_w$ from a Gnielinski or Dittus--Boelter correlation on the tube side, $h_a$ from a
finned-surface Colburn-$j$ correlation on the air side, and $\eta_o$ the overall fin
efficiency. Because $A_a\gg A_w$ the air-side term dominates --- another way of saying that the
physical size and cost of the coil live on the air side. Equation~\eqref{eq:hyd:UAdecomp} is
what one would use to design and price an actual coil; Eq.~\eqref{eq:hyd:UAcoil} is what the
system-level model needs, and the two should be reconciled before any hardware quotation is
taken seriously.
